For $\alpha\ne\mathcal O_P$, the standard vanishing for a nontrivial
line bundle in $\operatorname{Pic}^0(P)$ and the divisor sequence give
\[
 H^0(X,(L\otimes\alpha)|_X)\simeq H^0(P,L\otimes\alpha),
 \qquad h^0(P,L\otimes\alpha)=3.
\]
Combining this with \eqref{eq:canonical-resolution} and
$\rho_*\mathcal O_S(-E_{\exc})=\mathfrak m_0$ yields
\begin{equation}
 H^0(S,K_S\otimes j^*\alpha)
 \simeq
 \ker\!\left(
 H^0(P,L\otimes\alpha)\xrightarrow{\operatorname{ev}_0}
 (L\otimes\alpha)_0\right).
\label{eq:twisted-canonical-evaluation}
\end{equation}
The locus on which the evaluation map in
\eqref{eq:twisted-canonical-evaluation} is nonzero is a nonempty Zariski-open
subset of $\operatorname{Pic}^0(P)$.  For nonemptiness, use the
surjectivity of $\varphi_L$: every $L\otimes\alpha$ is a translate of $L$,
and a translate of the defining section by a point outside $X$ does not
vanish at the origin.  Intersecting this open with its inverse under
$\alpha\mapsto\alpha^{-1}$ gives a dense open $W\subset\operatorname{Pic}^0(P)$
on which
\begin{equation}
 h^0(S,K_S\otimes j^*\alpha)
 =h^0(S,K_S\otimes j^*\alpha^{-1})=2.
\label{eq:twisted-outer-hodge}
\end{equation}

Perverse generic vanishing gives a nonempty Zariski-open subset of the Betti
character torus on which
\[
 H^k(P,Rj_*\QQ_S[2]\otimes\mathcal L_\alpha)=0\qquad(k\ne0).
\]
The compact unitary character torus is Zariski dense in the Betti character
torus, and its unitary-character identification with
$\operatorname{Pic}^0(P)$ is a real-analytic isomorphism.  The restriction of
each of the two nonempty Zariski-open loci to the unitary torus is therefore
a dense Euclidean open subset: a proper complex algebraic subset cannot
contain a Euclidean open subset of this totally real Zariski-dense torus.
Their intersection is consequently a nonempty Euclidean open subset
$\Omega$ of the unitary character torus.

\begin{proposition}[Generic-twist Hodge vector]\label[proposition]{prop:hodge-vector}
For every $\alpha\in\Omega$, the Hodge structure
$H^0(P,\IC_X^H\otimes\mathcal L_\alpha)$ has Hodge numbers
\[
 (h^{2,0},h^{1,1},h^{0,2})=(2,3,2).
\]
Equivalently, the associated antidiagonal Hodge cocharacter has weights
$-2,0,2$ on $\omega(\IC_X)$ with multiplicities $2,3,2$, respectively.
\end{proposition}

\begin{proof}
For $\alpha\in\Omega$, perverse generic vanishing
\cite[Theorem~1.1]{KraemerWeissauer2015} gives
\[
 H^k(S,j^*\mathcal L_\alpha)=0\qquad(k\ne2).
\]
For a unitary rank-one local system, Hodge theory for the associated flat
unitary line bundle gives the twisted decomposition
\[
 H^{p,q}(S,j^*\mathcal L_\alpha)
 =H^q(S,\Omega_S^p\otimes j^*\alpha),
 \qquad p+q=2.
\]
Equation \eqref{eq:twisted-outer-hodge}, together with Serre duality, gives
outer Hodge numbers $2$ and $2$.  Since the Euler characteristic of a
rank-one local system equals $e(S)=8$ and all cohomology is concentrated in
degree two, the middle Hodge number is $4$.  Thus the Hodge vector of
$Rj_*\QQ_S^H[2]$ is $(2,4,2)$.  The Hodge-module refinement of
\eqref{eq:semismall}, extended to $P$, is
\[
 Rj_*\QQ_S^H[2]\simeq\IC_X^H\oplus\QQ_0^H(-1).
\]
The point summand has type $(1,1)$, so the Hodge vector for $\IC_X^H$ is
$(2,3,2)$.
\end{proof}

The symmetry of $X$ also gives the self-duality needed in the tensor
argument.  Indeed, Verdier self-duality of the intersection complex and
\eqref{eq:symmetry} give
\begin{equation}
 K^\vee=(-1_P)^*D(K)\simeq K.
\label{eq:self-duality}
\end{equation}
This is the self-duality used below; the same pairing reappears in Section~7.

\begin{lemma}[The chosen twist is tensor-compatible]
\label[lemma]{lem:tensor-compatible-twist}
For every $\alpha\in\Omega$, the functor
\[
 \omega_\alpha(Q)=H^0(P,Q\otimes\mathcal L_\alpha)
\]
on the semisimple tensor category generated by the image of $K=\IC_X$ in
$\overline{\operatorname{Perv}}(P)$ is a fibre functor.  In particular the
Hodge decomposition of \cref{prop:hodge-vector} defines the Hodge
cocharacter on the same Tannakian category.
\end{lemma}

\begin{proof}
The decomposition
$Rj_*\QQ_S[2]\simeq K\oplus\delta_0$ and the definition of $\Omega$ give
\[
 H^i(P,K\otimes\mathcal L_\alpha)=0\qquad(i\ne0).
\]
Let $m_n:P^n\to P$ be addition.  Since a rank-one character local system
satisfies
$m_n^*\mathcal L_\alpha\simeq\mathcal L_\alpha^{\boxtimes n}$, the
projection formula and K\"unneth give
\begin{equation}
 H^\bullet(P,K^{*n}\otimes\mathcal L_\alpha)
 \simeq
 H^\bullet(P,K\otimes\mathcal L_\alpha)^{\otimes n},
\label{eq:tensor-twist-kunneth}
\end{equation}
so the left-hand side is concentrated in degree zero.

The complexes $K^{*n}$ are projective direct images of pure polarizable
Hodge modules, hence the decomposition theorem splits each of their perverse
cohomology sheaves as direct summands.  Consequently
${}^pH^0(K^{*n})$, and every simple direct summand of it, has twisted
cohomology concentrated in degree zero.  All other perverse-cohomological
pieces are negligible after passage to
$\overline{\operatorname{Perv}}(P)$, because convolution is $t$-exact in
the quotient.  Moreover, any negligible simple summand of
${}^pH^0(K^{*n})$ has Euler characteristic zero; the preceding concentration
then forces its twisted $H^0$ to vanish.  Hence $H^0(-\otimes\mathcal
L_\alpha)$ factors through the quotient on the tensor category under
consideration.  Since $K^\vee\simeq K$, every object in the semisimple tensor
category generated by $K$ is a direct summand of a finite direct sum of such
${}^pH^0(K^{*n})$.  Thus the same vanishing holds for every object of that
category.

Equation \eqref{eq:tensor-twist-kunneth}, by naturality and passage to direct
summands, supplies the tensor constraint for $\omega_\alpha$; exactness and
faithfulness follow from semisimplicity and the equality
$\dim H^0(P,Q\otimes\mathcal L_\alpha)=\chi(Q)>0$ for every nonzero simple
object in the quotient.  This agrees with the generic-twist fibre-functor
formalism of \cite[Section~3.1, Proposition~3.1]{KLM2026}.
\end{proof}


\section{The generic pair-incidence splitting}\label{sec:pair-splitting}

Return to the family over $\cB$.  Let $\cE\to\cB$ be the Weierstrass elliptic
scheme with zero section $O$, put $\eta=\mathcal O_{\cE}(O)$, and let
\[
 \pi:\cC\longrightarrow\cE,
 \qquad
 \kappa=\pi^*\eta
\]
be the relative shifted cyclic cover and its theta characteristic.  Let
$\cP\to\cB$ be the relative Prym scheme and $\cX\subset\cP$ the relative
centred theta section.  For any $Y\to\cB$, undecorated symbols
$E_Y,C_Y,P_Y,X_Y,\eta$ and $\kappa$ denote the corresponding base changes.

The relative theta divisor determined by $\kappa$ restricts to a section of
a relatively ample line bundle on the abelian scheme $\cP\to\cB$.  Its
restriction to every closed fibre is the nonzero Cartier divisor constructed
in \cref{prop:theta-resolution}.  Hence its zero scheme
$\cX\subset\cP$ is a symmetric relative effective Cartier divisor; the
defining section is a nonzerodivisor on every fibre of the smooth morphism
$\cP\to\cB$, so $\cX\to\cB$ is flat.

